\begin{abstract}
We develop an operational isomorphism between \emph{complexity}, \emph{geometry}, and \emph{observation} within a unified syntax for Holographic Polar Arithmetic (HPA) and Omega Theory. In this framework, computational resources are not external bookkeeping but internal geometric costs of a scan--readout protocol: time is defined as the iteration count of a genuine unitary scan operator $\Theta$, while space is defined as readout resolution (prefix length and orthogonal cut capacity) induced by a window projection $\PiW$ and its canonical coding (Ostrowski numeration, specializing to Zeckendorf in the golden branch). The scan shift and phase multiplication form a Weyl pair, yielding a variance-type uncertainty relation that forbids simultaneous minimization of scan-localization and readout-localization, thereby providing a structural time--space complementarity.

On the dynamical side, Omega Theory models microscopic evolution as a partitioned quantum cellular automaton (PQCA) and supports an \emph{exact} compilation of a local PQCA step on a finite region into a one-dimensional nearest-neighbor circuit. The resulting compilation depth defines a routing overhead $\kappa$ and an emergent computational lapse field $\lapse(x)=\kappa_0/\kappa(x)$, making gravitational time dilation equivalent to computational slowdown in an operational sense (Appendix~\ref{app:compilation_overhead}). At the computability layer, universal QCA dynamics admits an undecidable local reachability predicate by reduction from reversible halting (Appendix~\ref{app:reachability_undecidability}), which provides a theorem-level ``open-endedness boundary'' for prediction.

We emphasize a strict separation of layers: undecidability is a theorem about internal predicates, whereas modeling observers as interactive machines with oracle-like interfaces is an \emph{interpretation-layer} language for resource injection and conditionalization, not a deduction about physical collapse. Finally, the discrete scan provides an exact lattice dispersion relation whose low-energy limit recovers approximately linear propagation and whose high-energy deviation yields a testable template for energy-dependent corrections to effective signal speed.
\end{abstract}

\textbf{Keywords}: Holographic Polar Arithmetic; Omega Theory; quantum cellular automata; time complexity; space complexity; Weyl pair; Ostrowski numeration; Zeckendorf decomposition; undecidability; interactive computation; oracle.

\paragraph{Conventions.}
Unless otherwise stated, $\log$ denotes the natural logarithm. ``mod $1$'' refers to reduction in $\RR/\ZZ$, and ``mod $2\pi$'' refers to reduction in $\RR/2\pi\ZZ$.

\tableofcontents
\newpage
